\chapter{Общий родитель длины ребра клеточного комплекса}
\label{ch:v10-nonequilibrium-bath-edge-length-common-parent}

Типизированная инцидентность ещё не задаёт метр. Проверим, могут ли все
полученные ранее относительные связи совместно выбрать физическую длину.
Обозначим через $\ell_{\rm edge}$ длину ребра причинного графа, а через
$\ell_{\rm cell}$ --- ребро изотропной четырёхмерной ячейки. Минимальный
единый пакет состоит из семи условий
\begin{align}
 \frac{v_{\rm cell}}{\ell_{\rm cell}^4}&=1,&
 \frac{E_C\ell_{\rm cell}}{\hbar c}&=1,&
 k_{\rm BZ}\ell_{\rm cell}&=\pi,\\
 \frac{\omega_{\rm UV}\ell_{\rm cell}}{v_g}&=2\sqrt3,&
 \Lambda_{43}\ell_{\rm cell}&=42,&
 \frac{v_g\Delta t}{\ell_{\rm edge}}&=1,\\
 \frac{\ell_{\rm edge}}{\ell_{\rm cell}}&=1.&&
\end{align}
Последнее равенство типизирует причинный путь как односкелетную часть того же
метрического комплекса. Из пакета следуют прежнее отношение
$\Lambda_{43}/k_{\rm BZ}=42/\pi$, равенство
$v_{\rm cell}=\ell_{\rm edge}^4$ и точная UV-фаза одного причинного шага
\begin{equation}
 \omega_{\rm UV}\Delta t=2\sqrt3.
\end{equation}

На семи безразмерных координатах положительный родитель, равный половине
суммы квадратов отклонений от
$(1,1,\pi,2\sqrt3,42,1,1)$, имеет гессиан $I_7$, ранг $7$ и определитель
$1$. Следовательно, противоречия между объёмным, часовым, спектральным и
причинным секторами нет.

Но в размерных переменных
\begin{equation}
 (v_{\rm cell},E_C,k_{\rm BZ},\omega_{\rm UV},\Lambda_{43},
 \Delta t,v_g,\ell_{\rm cell},\ell_{\rm edge})
\end{equation}
логарифмическая карта семи условий имеет ранг/ядро $7/2$. Одна мода меняет
единицу скорости. После физического отождествления $v_g=c$ остаётся мода
\begin{equation}
 (4,-1,-1,-1,-1,1,0,1,1),
\end{equation}
то есть одновременное растяжение обеих длин, четырёхобъёма и времени с
обратным сжатием всех энергий и волновых чисел. Только независимый якорь
$\ell_{\rm cell}$ повышает ранг до $9$.

\begin{theorem}[Совместимость относительных метрических данных]
Все известные объёмные, спектральные, часовые и причинные отношения допускают
единый строго положительный родитель. Он отождествляет длину ребра
распространения с ребром четырёхмерной ячейки и фиксирует все перечисленные
безразмерные отношения.
\end{theorem}

\begin{nogocriterion}[Общий родитель не является абсолютной линейкой]
После фиксации $v_g=c$ общий родитель сохраняет одномерную масштабную орбиту.
Поэтому ни собственный четырёхобъём, ни K43-обрезание, ни причинный шаг не
выбирают численное значение $\ell_{\rm edge}$ без независимого происхождения
глобального растущего комплекса или новой размерной границы.
\end{nogocriterion}

\begin{protocol}\begin{enumerate}
\item условная архитектура общего родителя: \textbf{$11/11$};
\item относительное отождествление двух рёбер: \textbf{$1/1$};
\item ранг/ядро размерной карты: \textbf{$7/2$};
\item после якоря $v_g=c$: \textbf{$8/1$};
\item происхождение глобального растущего комплекса: \textbf{$0/1$};
\item происхождение абсолютной длины ребра: \textbf{$0/1$};
\item следующий гейт: \textbf{типизированное вложение универсального
покрытия в граф рождения}.
\end{enumerate}\end{protocol}

Машинный сертификат сохранён в
\path{s2t_v10_cell_birth_four_volume_nonequilibrium_bath_cell_complex_edge_length_common_parent_origin_gate_results.json}.